% Einheit 3 von 5 – Punkte und Werte (bank/lineare-funktionen/e3.jsonl, zusammenbau v0.1)
%% TODO Zweigzeile Teil 1: was hier gelernt wird (Satz in Schülersprache aus dem Einheitstitel) – Einheit 3
\einheitenkopf[e3]{Einheit 3 von 5 · Punkte und Werte}
\zweigzeile{ab Kl. 8 · P10 oft}
\setcounter{aufgabe}{18}

%% TODO Titel: Ich-kann-Satz fehlt in der Bank (Platzhalter: Kettenname) – Nr. 19
%% TODO Anweisung: ein Satz über der ersten Teilaufgabe fehlt in der Bank – Nr. 19
\begin{aufgabe}{Punkt in Gleichung einsetzen}
\begin{teile}
\teil Einsetzen: Setze $x = 4$ in $3x - 2$ ein und rechne aus.
\end{teile}
\end{aufgabe}

%% TODO Titel: Ich-kann-Satz fehlt in der Bank (Platzhalter: Kettenname) – Nr. 20
%% TODO Anweisung: ein Satz über der ersten Teilaufgabe fehlt in der Bank – Nr. 20
\begin{aufgabe}{Funktionswert}
%% TODO Darstellung für schwach fehlt in der Bank (lineare-funktionen-e3-k2-s1-v1; Abschnitt „Für schwache Schüler“)
\swz{$f(4)$? Berechne den Funktionswert von $f(x) = 2x + 5$ an der Stelle $x = 4$.}{}
%% TODO Darstellung für schwach fehlt in der Bank (lineare-funktionen-e3-k2-s1-v2; Abschnitt „Für schwache Schüler“)
\swz{$f(5)$? Berechne den Funktionswert von $f(x) = 3x - 4$ an der Stelle $x = 5$.}{}
%% TODO Darstellung für schwach fehlt in der Bank (lineare-funktionen-e3-k2-s1-v3; Abschnitt „Für schwache Schüler“)
\swz{$f(3)$? Berechne den Funktionswert von $f(x) = -2x + 9$ an der Stelle $x = 3$.}{}
%% TODO Darstellung für schwach fehlt in der Bank (lineare-funktionen-e3-k2-s1-v4; Abschnitt „Für schwache Schüler“)
\swz{$f(2)$? Berechne den Funktionswert von $f(x) = 5x - 7$ an der Stelle $x = 2$.}{}
%% TODO Darstellung für schwach fehlt in der Bank (lineare-funktionen-e3-k2-s1-v5; Abschnitt „Für schwache Schüler“)
\swz{$f(2)$? Berechne den Funktionswert von $f(x) = -3x + 10$ an der Stelle $x = 2$.}{}
\swfrage{Was bleibt gleich, was ändert sich?}
%% TODO Darstellung für schwach fehlt in der Bank (lineare-funktionen-e3-k2-s2-v1; Abschnitt „Für schwache Schüler“)
\swz{$f(-3)$? Berechne den Funktionswert von $f(x) = 4x - 1$ an der Stelle $x = -3$.}{}
%% TODO Darstellung für schwach fehlt in der Bank (lineare-funktionen-e3-k2-s3-v1; Abschnitt „Für schwache Schüler“)
\swz{$f(1{,}5)$? Berechne den Funktionswert von $f(x) = 4x - 3$ an der Stelle $x = 1{,}5$.}{}
%% TODO Darstellung für schwach fehlt in der Bank (lineare-funktionen-e3-k2-s4-v1; Abschnitt „Für schwache Schüler“)
\swz{$\text{Welches x? Für welches x gilt $f(x) = 11$ bei $f(x) = 2x + 3$?}$}{}
%% TODO Darstellung für schwach fehlt in der Bank (lineare-funktionen-e3-k2-s5-v1; Abschnitt „Für schwache Schüler“)
\swz{Liegt der Punkt? Prüfe, ob $P(3|8)$ auf der Geraden $f(x) = 3x - 1$ liegt.}{}
%% TODO Darstellung für schwach fehlt in der Bank (lineare-funktionen-e3-k2-s6-v1; Abschnitt „Für schwache Schüler“)
\swz{$\text{Nullstelle? Berechne die Nullstelle von $f(x) = 3x - 12$.}$}{}
\swa{Graph und Nullstelle: Zeichne den Graphen der Funktion f mit $f(x) = -4x + 6$ und berechne die Nullstelle von f. \hfill (P10 2022 OS) x = \leerfeld}{\begin{ksys}[xmin=-4,xmax=4,ymin=-4,ymax=7] \end{ksys}}
\end{aufgabe}

%% TODO Titel: Ich-kann-Satz fehlt in der Bank (Platzhalter: Kettenname) – Nr. 21
%% TODO Anweisung: ein Satz über der ersten Teilaufgabe fehlt in der Bank – Nr. 21
\begin{aufgabe}{Wertetabelle}
\swa{Wertetabelle: Fülle die Tabelle zu $f(x) = -\frac{1}{2}x + 3$ aus.}{\wertetabelle{x}{f(x)}{-4,-2,0,2,4}}
\end{aufgabe}

%% TODO Titel: Ich-kann-Satz fehlt in der Bank (Platzhalter: Kettenname) – Nr. 22
%% TODO Anweisung: ein Satz über der ersten Teilaufgabe fehlt in der Bank – Nr. 22
\begin{aufgabe}{Schnittpunkt mit y-Achse}
%% TODO Darstellung für schwach fehlt in der Bank (lineare-funktionen-e3-k4-s1-v1; Abschnitt „Für schwache Schüler“)
\swz{Schnittpunkt mit der y-Achse? Gib ihn für $f(x) = 5x - 8$ an.}{}
\end{aufgabe}

%% TODO Titel: Ich-kann-Satz fehlt in der Bank (Platzhalter: Kettenname) – Nr. 23
%% TODO Anweisung: ein Satz über der ersten Teilaufgabe fehlt in der Bank – Nr. 23
\begin{aufgabe}{Funktionswert – Fehler finden, Begründen, Darstellungswechsel, Anwendung}
\swz[3]{Ich finde den Fehler: Sara soll die Nullstelle von $f(x) = 2x - 6$ angeben und schreibt: Die Nullstelle ist $-6$. Benenne den Fehler und rechne richtig.}{}
\swfrage{Zwei Nullstellen? Kann die Gerade $f(x) = 3x - 5$ zwei Nullstellen haben? Begründe.}
\swz{Gleichung aus der Tabelle: Gib die Gleichung der linearen Funktion f an.}{\wertetabelle[-1,2,5,8]{x}{f(x)}{0,1,2,3}}
\swz[3]{$\text{Wann ist der Aufzug unten? Ein Aufzug fährt aus 36 m Höhe gleichmäßig nach unten. Seine Höhe über dem Boden (in m) nach x Sekunden ist $f(x) = -1{,}5x + 36$. Nach wie vielen Sekunden ist er unten?}$}{}
\end{aufgabe}

\uebersichtskasten{Funktionswert: x einsetzen. \quad f(x) = 2x + 1, x = 3: f(3) = 2 · 3 + 1 = 7 \\ Punktprobe: x einsetzen, mit y vergleichen. \quad P(3 | 7) liegt auf f, weil f(3) = 7 (wA) \\ Nullstelle: f(x) = 0 setzen, nach x auflösen.  0 = 2x + 1 → x = −0,5}
